% !TeX root = ../main.tex
\section{Problem Definition}

Let the expected architecture contract be
\begin{equation}
  M = (C_E, E_E, R),
\end{equation}
where $C_E$ is the set of expected components, $E_E$ is the set of expected typed and directed relationships, and $R$ is the set of conformance rules. We denote the graph projection of this contract by $G_E=(C_E,E_E)$. A relationship is represented as a tuple $(u,t,v)$, where $u$ is the source component, $v$ is the target component, and $t$ is the relationship type.

Given a source-code repository $S$ and supplied model $M$, a stack-specific analyzer $A$ constructs an observed graph
\begin{equation}
  G_O = A(S;M) = (C_O, E_O, P),
\end{equation}
where $C_O$ and $E_O$ are observed components and relationships, and $P$ maps each observation to source evidence. The analyzer uses expected component identifiers to map source paths and reuses declared metadata for matching identifiers. This is model-conditioned reconstruction, not autonomous discovery of architectural boundaries. Evidence contains file, line, column, detector, excerpt, and a fixed confidence annotation, not a calibrated correctness probability.

ArchSync supports four deterministic rule semantics. A \emph{deny} rule rejects every matching relationship. An \emph{allow} rule constrains the outgoing relationships of selected sources to an allowed target set. A \emph{require} rule states that every matching source must have a direct edge to a matching target. A \emph{require-path} rule states that every matching source must reach a matching target through a non-empty directed path. Rules may optionally constrain the relationship type and may use wildcard component selectors.

The conformance function compares $M$ and $G_O$ and returns a graph difference $\Delta$, a set of rule violations $\mathcal{V}$, and a set of evidence-backed findings $\mathcal{F}$. Its classification is defined as
\begin{equation}
\operatorname{class}(M,G_O)=
\begin{cases}
\text{violation}, & \text{if } |\mathcal{V}| > 0,\\
\text{evolution}, & \text{if } |\mathcal{V}| = 0 \land |\Delta| > 0,\\
\text{no-impact}, & \text{otherwise.}
\end{cases}
\end{equation}

For a Git change from base repository $S_b$ to head repository $S_h$, ArchSync computes base and head finding sets and compares them by stable semantic identity:
\begin{equation}
\begin{aligned}
  \mathcal{F}_b &= \mathcal{F}(M,A(S_b;M)),\\
  \mathcal{F}_h &= \mathcal{F}(M,A(S_h;M)),\\
  K_b &= \{k(f):f\in\mathcal{F}_b\},\\
  \mathcal{F}^{+} &= \{f\in\mathcal{F}_h:k(f)\notin K_b\}.
\end{aligned}
\end{equation}
At the measured revision, $k(f)$ comprises kind, rule identifier, change, component, and canonical edge key; absent fields use empty strings. Source coordinates, messages, severity, and generated identifiers are not key fields. Thus an additional call site for an old violating edge need not introduce a finding. This is a relation-level policy, not an occurrence-level rejection policy.

Only $\mathcal{F}^{+}$ controls the pull-request decision. A new rule finding yields BLOCK; otherwise new evolution findings yield REVIEW; otherwise the result is PASS. Head findings with base keys are pre-existing; disappeared base keys are resolved. PASS certifies neither a violation-free head nor complete observation of source interactions.

The feasibility question is whether this model-conditioned analyzer and relation-level decision policy agree with declared controlled oracles. That agreement is a prerequisite for, not evidence of, effectiveness on independently labeled repository histories.

\section{Research Questions}

\textbf{RQ1: To what extent do model-conditioned graphs and detector signals agree with the declared D1/D2 annotations?}

For the 20-patch development benchmark (D1), this question reports matched, extra, and missed graph elements and patch deltas. For the separate 40-signal corpus (D2), it reports descriptive precision, recall, F1, true negatives, and specificity against the annotated detector signals; these are controlled-corpus results, not population estimates.

\textbf{RQ2: To what extent do classifications and violated-rule sets agree with declared controlled-case outcomes?}

This question evaluates exact classification agreement with a manually curated ground truth and, for violations, agreement on the expected rule identifiers.

\textbf{RQ3: To what extent do call-site and missing-relation anchor locations agree with declared file/line references at the case level?}

This question separates localization of observed calls from deterministic anchors for missing relationships. It measures case-level file/line agreement, not recall of all findings or call sites. Column, detector, excerpt, and model path remain diagnostic evidence.

\textbf{RQ4: Are ArchSync full-repository and Git-diff outputs deterministic and reproducible, and what analysis scope and latency are observed for the cached incremental gate?}

This question compares normalized observed graphs, findings, and pull-request decisions across repeated executions and clean installations. For Git-diff checks it additionally records cache behavior, files parsed relative to the head repository, and cold/warm latency on the measured environment.

These scope-specific RQ formulations clarify reporting after the controlled runs; original inputs, labels, metrics, and acceptance criteria are unchanged.

\section{Proposed Approach}

\subsection{Architecture as an Executable Contract}

ArchSync uses a YAML document as the source of truth for architecture intent. The model defines stable component identifiers, component types, layers, directed relationships, rule severity, ownership metadata, and optional quality goals. JSON Schema validation rejects malformed documents, while semantic validation rejects unknown component references, duplicate edges, invalid identifiers, and other topology inconsistencies. Mermaid and draw.io views are derived from the model; they are not independent sources of truth.

Stable identifiers and canonical edge keys normalize the validated graph for shared command-line, benchmark and CI conformance semantics.

\subsection{Source-to-Graph Analysis}

A Guardian analyzer traverses TypeScript files under component roots and parses them through the TypeScript Compiler API. Rather than matching arbitrary text, detectors inspect syntax-tree nodes, imports, call expressions, variable initialization, environment-variable references, and URL literals. Stack-specific detectors translate these signals into typed architectural relationships.

The detector set recognizes HTTP requests, PostgreSQL clients and queries, Redis clients and operations, and AMQP publishing and consumption. Version 0.2 tracks local bindings introduced by named, aliased, and namespace imports from \texttt{pg}, \texttt{redis}, and \texttt{amqplib}. PostgreSQL and Redis edges require an operation on an identifier associated with a recognized client initializer; AMQP edges require a channel associated with a tracked connection. These associations use per-file identifier names rather than complete compiler symbol or alias resolution. This package provenance filters same-named methods on unrelated objects when their receivers lack a tracked association; it does not establish soundness under arbitrary shadowing or reassignment. Endpoint resolution uses explicit literals and deterministic environment-variable fallbacks when available. Every detected relationship carries one or more source locations. Detected components and relationships are sorted and deduplicated before serialization.

\subsection{Conformance and Decision Semantics}

Source paths map to the longest expected component identifier that prefixes the relative path, falling back to the first path segment. URL targets use the lowercase hostname, ignoring port and path. Unresolved HTTP endpoints produce no edge and no explicit unknown finding for that call. Consequently, a deny-only contract can pass despite an unobserved interaction. AMQP publishing produces a producer-to-broker edge and consumption a broker-to-consumer edge. A typed require-path rule requires every traversed edge to have that type; an untyped rule permits mixed types. Broker-level reachability does not establish matching queues, routing keys, or message delivery.

The conformance engine first computes node and edge differences between the expected and observed graphs. It then evaluates the observed graph against the architecture rules. A rule violation takes precedence over ordinary topology evolution. This is a fixed conservative policy: one rule violation determines the top-level classification even when the change also contains acceptable evolution. Both types of findings remain available for inspection. Rule severity is reported but does not alter this precedence or the blocking threshold. This may delay acceptable changes bundled with a violation; alternative policies are not evaluated.

Full-repository checks map no-impact, violation and evolution to PASS, BLOCK and REVIEW, respectively; Git-diff checks apply this precedence only to introduced findings.
The REVIEW outcome separates a request for a design decision from a detected rule breach. In the evaluated command, PASS, BLOCK, and REVIEW return exit codes 0, 1, and 3, respectively. A repository must configure its merge controls to treat the nonzero REVIEW result as requiring resolution; an annotation alone does not enforce approval. ArchSync reports the change without automatically modifying the approved architecture contract.

\subsection{Diff-Scoped Living Architecture Loop}

The Phase 3 gate resolves the merge base, enumerates modified, added, deleted, renamed, and untracked files, and maps changed TypeScript paths to source components. A cold check reconstructs the base Observed Graph from the committed source tree. The graph is cached inside Git metadata under a key derived from the base commit, architecture-contract hash, and analyzer version. A warm check validates and reuses that cache.

For the proposed head, Guardian parses all TypeScript files in affected source components, rather than only changed lines, retaining independently detectable relationships elsewhere in those components. This scope does not provide inter-file client/data-flow resolution. Those components' observations replace their baseline counterparts; unaffected components remain cached. The same supplied model is used for analysis and conformance at base and head; separate historical models are not loaded. Editing the model invalidates the cache and changes the contract used for both snapshots, not a model-to-model finding delta. Cache-envelope checks assume trusted local contents and a maintained analyzer version; they do not authenticate the graph. Annotations, a report, and exit codes expose the decision.

The analyzer updates observations, never the model. Owner-reviewed model changes are a proposed acceptance workflow, not authenticated approvals in the evaluated gate. It does not persist accepted/rejected review decisions. If evolution is rejected but code and contract stay unchanged, the next run returns REVIEW again, not BLOCK. Repository controls must prevent merging; the team revises code or contract and reruns. Rationale recording and new rules for lasting constraints remain manual governance.

\subsection{Evidence-Rich Findings}

A finding contains a contract version, identifier, kind, severity, message, optional rule, affected element, source evidence, and model evidence. A forbidden edge retains recognized supporting call sites, not necessarily a newly added occurrence. A missing required edge or path uses a deterministic component anchor and the model rule: an absent call has no call site to localize.

\textbf{End-to-end example.} Existing D1 case 06 adds a frontend \texttt{fetch} with a payment-service URL fallback at \texttt{frontend/src/app.ts:14}. Source and hostname mapping yield $(\texttt{frontend},\mathrm{http},\texttt{payment\mbox{-}service})$. The edge matches deny rule ARCH-001. Its key is absent at the base, so the linked rule and call evidence yield BLOCK even though an evolution finding is also present. This is an existing controlled case, not an additional evaluation sample.

\begin{table*}[t]
\centering\small
\caption{Relation-level gate boundaries. Outcomes depend on supported observations; they do not guarantee complete detection.}
\label{tab:policy-boundaries}
\begin{tabular}{p{0.26\textwidth}p{0.28\textwidth}p{0.36\textwidth}}
\toprule
Condition & Gate behavior & What remains outside the guarantee \\
\midrule
New observed rule-finding key & BLOCK & Unsupported or unresolved interactions. \\
New evolution, no new rule finding & REVIEW & Human approval or automatic rejection handling. \\
No new key on a drifting base & PASS; old findings remain visible & A violation-free head. \\
New call site for an old violating edge & Can PASS; key unchanged & Rejection of every violating occurrence. \\
Unresolved HTTP target & No edge or explicit unknown finding & Complete observation of communication. \\
Code and model edited together & Same model evaluates both snapshots & Authorization of the changed contract. \\
\bottomrule
\end{tabular}
\end{table*}
